% Terjemahan Bahasa Indonesia dari chapter1.tex baris 1141--1229.
% Sumber: Methods of Algebra, Volume 2, commit
% 9a5803ff77dd3257484cb177f851a73770a59dd3 (CC BY 4.0).
% stable-unit-id: o014.aljabr2.chapter1.kan-extensions-as-limits

% segment-id: o014.li2.chapter1.kan-extensions-as-limits.h001
\section{Konstruksi Ekstensi Kan melalui Limit}\label{sec:Kan-as-limit}
\hypertarget{o014.aljabr2.chapter1.kan-extensions-as-limits}{ }

% segment-id: o014.li2.chapter1.kan-extensions-as-limits.p001
Sebagaimana dalam \S\ref{sec:Kan-extension}, pada bagian ini kita tetap
meninjau funktor $K: \mathcal{C} \to \mathcal{D}$ dan
$F: \mathcal{C} \to \mathcal{E}$, dengan tujuan mengkaji keberadaan
$\Lan_K F$ dan $\Ran_K F$.

% segment-id: o014.li2.chapter1.kan-extensions-as-limits.p002
Ingat kembali pembahasan pada awal \S\ref{sec:Kan-extension}. Ekstensi Kan
dimotivasi oleh upaya mencari funktor $G: \mathcal{D} \to \mathcal{E}$
sedemikian sehingga $F \simeq GK$, atau setidaknya hampiran terbaiknya.
Bagaimana kita mendefinisikan $Gd$ untuk $d \in \Obj(\mathcal{D})$?
Satu-satunya petunjuk ialah bahwa ketika $d = Kc$, objek $Gd$ seharusnya
diambil sebagai $Fc$ hingga isomorfisme. Hal ini mengilhami kita untuk
menghampiri $Gd$ dengan $\varinjlim$ (atau $\varprojlim$) dari semua $Fc$,
dengan limit diambil atas semua $c \in \Obj(\mathcal{C})$ dan morfisme
$Kc \to d$ (atau $d \to Kc$). Kedua arah limit ini akan memiliki sifat
universal yang saling dual. Berikut ini kita uraikan konstruksi tersebut.

% segment-id: o014.li2.chapter1.kan-extensions-as-limits.p003
Diberikan $d \in \Obj(\mathcal{D})$, definisikan kategori $(K/d)$ dan $(d/K)$
seperti dalam Definisi \ref{def:comma-category}, beserta funktor proyeksi
% segment-id: o014.li2.chapter1.kan-extensions-as-limits.q001
\[ \Pi_{/d}: (K/d) \to \mathcal{C}, \quad \Pi_{d/}: (d/K) \to \mathcal{C}. \]
Selanjutnya perhatian kita tertuju pada funktor komposit
$F \Pi_{/d}: (K/d) \to \mathcal{E}$ dan
$F \Pi_{d/}: (d/K) \to \mathcal{E}$, yang masing-masing memetakan objek
$(c, Kc \to d)$ dan $(c, d \to Kc)$ ke $Fc$. Dengan mengasumsikan bahwa
limit-limit yang dibicarakan ada, mula-mula kita catat beberapa pengamatan.

% segment-id: o014.li2.chapter1.kan-extensions-as-limits.l001
\begin{itemize}
	% segment-id: o014.li2.chapter1.kan-extensions-as-limits.i001
	\item Setiap morfisme $d \to d'$ menginduksi morfisme kanonik
	$\varinjlim F \Pi_{/d} \to \varinjlim F\Pi_{/d'}$; morfisme ini dicirikan
	oleh kekomutatifan diagram
	% segment-id: o014.li2.chapter1.kan-extensions-as-limits.q002
	\begin{equation}\label{eqn:Kan-limit-aux-0}
	% segment-id: o014.li2.chapter1.kan-extensions-as-limits.g001
	\begin{tikzcd}
		Fc \arrow[equal, r] \arrow[d] & Fc \arrow[d] \\
		\varinjlim F \Pi_{/d} \arrow[r] & \varinjlim F \Pi_{/d'}
	\end{tikzcd}\end{equation}
	untuk setiap $(c, Kc \to d) \in \Obj((K/d))$ (yang menginduksi
	$(c, Kc \to d') \in \Obj((K/d'))$), sedangkan anak panah vertikal adalah
	morfisme struktur kanonik untuk $\varinjlim$. Hal ini merupakan perwujudan funktorialitas
	limit; lihat \cite[Lema 2.7.4]{Li1}\footnote{Catatan penerjemah: sumber
	menulis lokator Lema~2.7,4; target memperbaikinya menjadi Lema~2.7.4
	(O014-C016).} beserta buktinya.
	% segment-id: o014.li2.chapter1.kan-extensions-as-limits.i002
	\item Dengan meninjau
	$(c, Kc \xrightarrow{\identity} Kc) \in \Obj(K/Kc)$, kita memperoleh
	morfisme kanonik
	$\eta_c: Fc \to \varinjlim F \Pi_{/Kc}$.
	% segment-id: o014.li2.chapter1.kan-extensions-as-limits.i003
	\item Setiap morfisme $c \to c'$ menginduksi morfisme kanonik
	$\varinjlim F \Pi_{/Kc} \to \varinjlim F \Pi_{/Kc'}$ yang membuat diagram
	berikut komutatif:
	% segment-id: o014.li2.chapter1.kan-extensions-as-limits.g002
	\[\begin{tikzcd}
		Fc \arrow[r] \arrow[d, "{\eta_c}"'] & Fc' \arrow[d, "{\eta_{c'}}"] \\
		\varinjlim F \Pi_{/Kc} \arrow[r] & \varinjlim F \Pi_{/Kc'} .
	\end{tikzcd}\]
\end{itemize}

% segment-id: o014.li2.chapter1.kan-extensions-as-limits.p004
Berdasarkan dualitas, $\varprojlim F \Pi_{d/}$ tentu memiliki sifat-sifat yang
bersesuaian, termasuk morfisme kanonik
$\varepsilon_c: \varprojlim F\Pi_{Kc/} \to Fc$. Pernyataan teorema berikut
akan didasarkan pada funktorialitas ini.

% segment-id: o014.li2.chapter1.kan-extensions-as-limits.th001
\begin{teorema}\label{prop:Kan-limit}
	Tinjau funktor $K: \mathcal{C} \to \mathcal{D}$ dan
	$F: \mathcal{C} \to \mathcal{E}$.
	% segment-id: o014.li2.chapter1.kan-extensions-as-limits.l002
	\begin{enumerate}[(i)]
		% segment-id: o014.li2.chapter1.kan-extensions-as-limits.i004
		\item Jika untuk setiap $d \in \Obj(\mathcal{D})$, limit
		$\varinjlim (F \Pi_{/d})$ ada di $\mathcal{E}$, maka
		$(\Lan_K F)(d) := \varinjlim (F \Pi_{/d})$, bersama dengan
		funktorialitas yang diberikan oleh \eqref{eqn:Kan-limit-aux-0},
		menentukan ekstensi Kan kiri
		$\Lan_K F: \mathcal{D} \to \mathcal{E}$. Transformasi alami yang
		bersesuaian, $\eta: F \to (\Lan_K F) \circ K$, berasal dari
		$(\eta_c)_{c \in \Obj(\mathcal{C})}$ yang dibahas sebelumnya.
		% segment-id: o014.li2.chapter1.kan-extensions-as-limits.i005
		\item Jika untuk setiap $d \in \Obj(\mathcal{D})$, limit
		$\varprojlim (F \Pi_{d/})$ ada di $\mathcal{E}$, maka
		$(\Ran_K F)(d) := \varprojlim (F \Pi_{d/})$, bersama dengan versi dual
		dari \eqref{eqn:Kan-limit-aux-0}, menentukan ekstensi Kan kanan
		$\Ran_K F: \mathcal{D} \to \mathcal{E}$. Transformasi alami yang
		bersesuaian, $\varepsilon: (\Ran_K F) \circ K \to F$, berasal dari
		$(\varepsilon_c)_{c \in \Obj(\mathcal{C})}$ yang dibahas sebelumnya.
		% segment-id: o014.li2.chapter1.kan-extensions-as-limits.i006
		\item Jika $K: \mathcal{C} \to \mathcal{D}$ penuh dan setia, maka
		$\eta$ yang dikonstruksi dalam (i) (atau $\varepsilon$ yang dikonstruksi
		dalam (ii)) merupakan isomorfisme.
	\end{enumerate}
\end{teorema}
% segment-id: o014.li2.chapter1.kan-extensions-as-limits.pf001
\begin{bukti}
	% segment-id: o014.li2.chapter1.kan-extensions-as-limits.p005
	Berdasarkan dualitas, berikut ini kita hanya membahas kasus $\Lan_K F$.
	Untuk (i), kita akan memeriksa satu per satu sifat universal dalam
	Definisi \ref{def:Kan-extension}.

	% segment-id: o014.li2.chapter1.kan-extensions-as-limits.p006
	Diberikan funktor $L: \mathcal{D} \to \mathcal{E}$ dan transformasi alami
	$\xi: F \to LK$. Untuk setiap $d \in \Obj(\mathcal{D})$, setiap objek
	$(c, Kc \to d)$ dalam $(K/d)$, serta setiap morfisme dalam $(K/d)$ yang
	ditentukan oleh $f: c \to c'$,
	$(c, Kc \to d) \to (c', Kc' \to d)$, terdapat diagram komutatif
	% segment-id: o014.li2.chapter1.kan-extensions-as-limits.g003
	\[\begin{tikzcd}
		Fc \arrow[r, "\xi_c"] \arrow[d, "Ff"'] & LKc \arrow[d, "LKf"'] \arrow[r] & Ld . \\
		Fc' \arrow[r, "{\xi_{c'}}"'] & LKc' \arrow[ru] &
	\end{tikzcd}\]
	Dengan menerapkan sifat universal $\varinjlim$, kita memperoleh
	$\chi_d: (\Lan_K F)(d) \to Ld$, yang dicirikan oleh kekomutatifan diagram
	% segment-id: o014.li2.chapter1.kan-extensions-as-limits.q003
	\begin{equation}\label{eqn:Kan-limit-aux-1}
	% segment-id: o014.li2.chapter1.kan-extensions-as-limits.g004
	\begin{tikzcd}
		Fc \arrow[r] \arrow[d, "\xi_c"'] & (\Lan_K F)(d) \arrow[d, "\chi_d"] \\
		LKc \arrow[r, "{L(Kc \to d)}"'] & Ld
	\end{tikzcd}\end{equation}
	untuk semua $(c, Kc \to d) \in \Obj((K/d))$; anak panah pada baris pertama
	merupakan morfisme struktur kanonik untuk $\varinjlim$.

	% segment-id: o014.li2.chapter1.kan-extensions-as-limits.p007
	Sekarang kita buktikan bahwa $(\chi_d)_{d \in \Obj(\mathcal{D})}$
	memberikan transformasi alami $\chi: \Lan_K F \to L$. Untuk suatu morfisme
	$d \to d'$, pilih objek $(c, Kc \to d)$ beserta citranya
	$(c, Kc \to d')$. Dengan membandingkan
	\eqref{eqn:Kan-limit-aux-0} dan \eqref{eqn:Kan-limit-aux-1}, persoalannya
	tereduksi menjadi membuktikan kekomutatifan diagram
	% segment-id: o014.li2.chapter1.kan-extensions-as-limits.g005
	\[\begin{tikzcd}[row sep=tiny]
		& & Ld \arrow[dd, "{L(d \to d')}"] \\
		Fc \arrow[r, "\xi_c"] & LKc \arrow[ru, "{L(Kc \to d)}"] \arrow[rd, "{L(Kc \to d')}"'] & \\
		& & Ld'
	\end{tikzcd}\]
	yang jelas berlaku.

	% segment-id: o014.li2.chapter1.kan-extensions-as-limits.p008
	Langkah berikutnya adalah memeriksa
	% segment-id: o014.li2.chapter1.kan-extensions-as-limits.q004
	\begin{equation}\label{eqn:Kan-limit-aux-2}
		\forall c \in \Obj(\mathcal{C}), \quad \xi_c = \chi_{Kc} \eta_c: Fc \to LKc.
	\end{equation}
	Karena $\eta_c$ merupakan morfisme kanonik dari
	$Fc = F\Pi_{/Kc}(c, Kc \xrightarrow{\identity} Kc)$ ke
	$(\Lan_K F)(Kc)$, persamaan \eqref{eqn:Kan-limit-aux-2} tidak lain adalah
	kasus khusus \eqref{eqn:Kan-limit-aux-1}.

	% segment-id: o014.li2.chapter1.kan-extensions-as-limits.p009
	Terakhir, kita tunjukkan bahwa transformasi alami
	$\chi: \Lan_K F \to L$ yang memenuhi \eqref{eqn:Kan-limit-aux-2} bersifat
	unik. Tetapkan $d$ dan tinjau sembarang objek $(c, Kc \to d)$ dalam $(K/d)$
	beserta diagram terkait
	% segment-id: o014.li2.chapter1.kan-extensions-as-limits.g006
	\[\begin{tikzcd}
		& Fc \arrow[ld, "{\eta_c}"'] \arrow[rd] & \\
		(\Lan_K F)(Kc) \arrow[rr, "{(\Lan_K F)(Kc \to d)}"'] \arrow[d, "{\chi_{Kc}}"'] & & (\Lan_K F)(d) \arrow[d, "\chi_d"] \\
 		LKc \arrow[rr, "{L(Kc \to d)}"'] & & Ld
	\end{tikzcd}\]
	di mana $Fc \to (\Lan_K F)(d)$ adalah morfisme yang muncul dalam
	\eqref{eqn:Kan-limit-aux-1}. Persegi itu komutatif karena kealamian
	$\chi$, sedangkan segitiganya komutatif berdasarkan
	\eqref{eqn:Kan-limit-aux-0}. Karena
	$\chi_{Kc} \eta_c = \xi_c$, komposisi
	$Fc \to (\Lan_K F)(d) \xrightarrow{\chi_d} Ld$ sama dengan
	$L(Kc \to d) \xi_c$. Dengan memvariasikan $(c, Kc \to d)$, keluarga
	persamaan ini menentukan $\chi_d$ secara unik.

	% segment-id: o014.li2.chapter1.kan-extensions-as-limits.p010
	Dengan demikian, $\left( \Lan_K F, \eta \right)$ memang merupakan ekstensi
	Kan kiri.

	% segment-id: o014.li2.chapter1.kan-extensions-as-limits.p011
	Sekarang tinjau (iii). Tetapkan $c \in \Obj(\mathcal{C})$. Karena $K$ penuh
	dan setia, mudah dilihat bahwa $(K/Kc)$ mempunyai objek terminal
	$(c, Kc \xrightarrow{\identity} Kc)$. Morfisme
	$\eta_c: Fc \to \varinjlim F\Pi_{/Kc}$ berasal dari objek terminal ini dan
	oleh karena itu merupakan isomorfisme.
\end{bukti}

% segment-id: o014.li2.chapter1.kan-extensions-as-limits.r001
\begin{catatan}
	Jika $\mathcal{C}$ merupakan kategori kecil, maka $(K/d)$ dan $(d/K)$ juga
	merupakan kategori kecil untuk setiap $d \in \Obj(\mathcal{D})$. Oleh karena
	itu, jika $\mathcal{E}$ kolengkap (atau lengkap), Teorema
	\ref{prop:Kan-limit} dan Proposisi \ref{prop:Kan-uniqueness} menunjukkan
	bahwa $K^*: \mathcal{E}^{\mathcal{D}} \to
	\mathcal{E}^{\mathcal{C}}$ mempunyai funktor adjoin kiri $\Lan_K$ (atau
	funktor adjoin kanan $\Ran_K$).
\end{catatan}

% segment-id: o014.li2.chapter1.kan-extensions-as-limits.e001
\begin{contoh}[Citra invers praberkas (presheaf)]
	\index{praberkas (presheaf)}
	Ambil $\mathcal{E} := \cate{Set}$, yang merupakan kategori lengkap dan
	kolengkap. Tinjau pemetaan kontinu $f: X \to Y$ di antara ruang-ruang
	topologis. Tetapkan $\mathcal{C} := \cate{Open}_Y$ sebagai kategori yang
	objeknya adalah subhimpunan terbuka di $Y$ dan morfismenya adalah pemetaan
	inklusi di antara subhimpunan terbuka; secara serupa, tetapkan
	$\mathcal{D} := \cate{Open}_X$. Sekarang definisikan funktor
	$\cate{Open}_Y \to \cate{Open}_X$ yang memetakan subhimpunan terbuka
	$V \subset Y$ ke $f^{-1} V$, dan selanjutnya pandang funktor ini sebagai
	$K: \cate{Open}_Y^{\opp} \to \cate{Open}_X^{\opp}$. Pembaca yang mengenal
	teori berkas akan segera melihat bahwa objek dalam
	$\cate{Open}_Y^\wedge := \cate{Set}^{\cate{Open}_Y^{\opp}}$, menurut
	definisi, tidak lain adalah praberkas pada $Y$, sedangkan
	$K^*: \cate{Open}_X^\wedge \to \cate{Open}_Y^\wedge$ merupakan funktor
	citra langsung di antara kategori praberkas, yang memetakan praberkas
	$\mathcal{F}$ pada $X$ ke praberkas pada $Y$ yang diberikan oleh
	$V \mapsto \mathcal{F}(f^{-1} V)$. Dalam situasi ini, adjoin kiri dari
	citra langsung, yaitu $\Lan_K$, diberikan oleh Teorema
	\ref{prop:Kan-limit}; funktor ini memetakan praberkas $\mathcal{G}$ pada
	$Y$ ke
	% segment-id: o014.li2.chapter1.kan-extensions-as-limits.q005
	\begin{gather*}
		U \mapsto \varinjlim_{\substack{V \subset Y: \text{subhimpunan terbuka} \\ \text{yang memenuhi}\; U \subset f^{-1} V }} \mathcal{G}(V), \quad U \subset X: \text{subhimpunan terbuka}, \\
		U \subset f^{-1}V \;\xleftrightarrow{\text{bersesuaian dengan}}\; \text{morfisme}\; K(V) \to U \;\text{di}\; \cate{Open}_X^{\opp}.
	\end{gather*}
	Tidak mengherankan, inilah tepatnya \emph{citra invers} dari
	$\mathcal{G}$ dalam teori berkas.
\end{contoh}
